\label{sec:introduction}

% \vspace{-4ex}
\section{Introduction}

Multimodal Large Language Models (MLLMs) \cite{alayrac2022flamingo, liu2023llava, dai2023instructblip, wang2023cogvlm, liu2023improvedllava, chen2023internvl, dong2024internlmxcomposer24khd, chen2024far, laurençon2023obelics, laurenccon2024matters, bai2023qwenvl, openai2024gpt4, claude3, rekateam2024reka, team2023gemini, li2024minigemini} are highly versatile and effective for a wide range of real-world applications \cite{masry2022chartqa, mathew2021docvqa, fu2024mme, liu2024mmbench, yu2023mm, lu2022learn, mathew2022infographicvqa, lu2024mathvista, yue2023mmmu}. 
Within these applications, chart understanding is a much desired capability as charts are ubiquitous in scientific papers, financial reports, and news articles. It also poses unique challenges where models need to perform complex reasoning over numerical data, textual labels, and complex visual elements to answer difficult questions (see \cref{fig:teaser}), thus making chart understanding a suitable measure of progress for MLLMs.
Many benchmarks in the popular MathVista evaluation suite~\citep{lu2024mathvista} are designed to test chart understanding.
However, these benchmarks lack diversity in both the types and complexity of the charts and the often template-based questions (\cref{subsec:design_flaw}). For example, FigureQA~\citep{kahou2017figureqa} and DVQA~\citep{kafle2018dvqa} rely on procedurally generated question templates. While ChartQA~\citep{masry2022chartqa} includes a mixture of handwritten and machine-generated questions, the charts lack visual diversity due to the homogeneous appearance of the charts from a limited number of sources. 
Regardless, many proprietary models \cite{openai2024gpt4, team2023gemini, claude3, rekateam2024reka} and open-source models \cite{chen2024far,dong2024internlmxcomposer24khd, chen2023internvl,laurenccon2024matters,viscpm,li2024minigemini,liu2024llavanext,gao2024sphinxx} evaluate on these datasets.\footnote{We note that there are several more sophisticated benchmarks \cite{xu2024chartbench, xia2024chartx, liu2024mmc} that have recently been released. We discuss key differences between CharXiv and these benchmarks in \cref{subsec:design_flaw}.}
These narrow evaluations have given the appearance that the open-source models outperform proprietary ones\footnote{See the FQA (\textit{i.e.,} Figure QA) column of the MathVista \href{https://mathvista.github.io/}{leaderboard}. Throughout the paper, ``open-source'' refers to models with publicly available weights.}, despite evidence to the contrary: 
we design simple stress tests (\cref{subsec:control}) in which we find that open-source models lag far behind proprietary ones in their robustness to small visual or textual changes. For example, the accuracy of SPHINX V2 dropped from 63.2\% to 28.6\% with a 34.5\% gap when questions are slightly modified with respect to the same set of charts. 
\begin{figure*}[t!]
  \centering
    \includegraphics[width=1\linewidth]{figures/teaser_new.pdf}
    \caption{Example chart (left), descriptive questions (top-right) and reasoning questions (bottom-right) in \ours{} where open-source models even fail in basic descriptive questions. Moreover, all models struggle with correctly answering the reasoning question.}
    \vspace{-4ex}
    \label{fig:teaser}
\end{figure*}

We introduce \ours{}, a comprehensive evaluation suite for complex understanding of natural, challenging, and diverse charts (\cref{sec:charxiv_overview}) to address the above issue.
\ours{} consists of 2,323 real-world charts handpicked from scientific papers spanning 8 major subjects published on arXiv (\cref{subsec:charxiv_chart}). 
We explicitly disentangle visual understanding and reasoning by designing two types of questions (\cref{subsec:charxiv_questions}): (1) \emph{descriptive} questions, requiring understanding basic chart information such as the title, labels, and ticks; (2) \emph{reasoning} questions, requiring comparisons, approximations, and fine-grained analysis.
\ours{} is an especially high-quality dataset where all questions are \textit{manually} curated by human experts, and all ground-truth answers are validated by hand.
To answer both types of questions, the model only needs to understand the visual contents of the chart without advanced domain-specific knowledge and contextual information.
Evaluating an MLLM on \ours{} is straightforward, because we impose a short answer format that is amenable to LLM-based automatic grading.

We extensively evaluate $13$ open-source models and $11$ proprietary models (\cref{subsec:exp_setup}) and identify a large disparity between the strongest open-source and proprietary models (\cref{subsec:exp_results}): InternVL Chat V1.5 correctly answers only $29.2$\% of the reasoning questions and $58.5$\% of the descriptive ones, whereas GPT-4o correctly answers $47.1$\% of the reasoning questions and $84.5$\% of the descriptive ones (\cref{tab:main_result}).
As shown in \cref{fig:benchmark_comparison}, the performance gap in the reasoning questions of $17.9\%$ is significantly larger than the gap identified in prior works \cite{kafle2018dvqa, kahou2017figureqa, masry2022chartqa}.
Further, both types of models lag far behind the human performance of $80.5$\% on the reasoning questions and $92.1$\% on the descriptive ones.
Fine-grained analysis of model performance (\cref{subsec:exp_analysis}) shows several insights owing to the design of \ours{}. 
In particular, we characterize: (1) differences in reasoning and descriptive capabilities, exploring when one skill reinforces the other; (2) what types of tasks and charts are difficult for existing MLLMs; (3) how different MLLMs respond to unanswerable questions.
Overall, we hope that \ours{} enables a thorough, multi-faceted evaluation of chart understanding in MLLMs. 

\begin{figure*}[t]
  \centering
    \vspace{-1ex}
    \includegraphics[width=1\linewidth]{figures/comparison.pdf}
    \vspace{-3ex}
    \caption{Model performance comparison on reasoning questions from \ours{} v.s. questions from existing benchmarks.
    As indicated by the red and blue bars resepctively, many open-source models surpass proprietary model performance on the $174$ sample questions from existing benchmarks (subsets of DVQA, FigureQA and ChartQA from the \textit{testmini} split of MathVista) yet fail consistently on the $1000$ reasoning questions from the validation split of \ours{}.}
    \label{fig:benchmark_comparison}
    \vspace{-3ex}
\end{figure*}
